% !TEX root = ../main.tex
% ══════════════════════════════════════════════════════════════
\chapter{Choix technologiques et algorithmiques}
\label{cha:choix-technologiques}
% ══════════════════════════════════════════════════════════════
\begingroup
\let\clearpage\relax
\startcontents[chapters]
\printcontents[chapters]{}{1}{\section*{Plan}}
\endgroup
\newpage

\section*{Introduction}
\justifying
\phantomsection
\addcontentsline{toc}{section}{Introduction}

Le chapitre précédent a montré qu'aucune solution existante ne répondait aux besoins du service RI, d'où le choix de développer une plateforme sur mesure. Ce chapitre justifie les technologies et les algorithmes retenus pour la construire, en comparant à chaque fois les principales alternatives.

% ─────────────────────────────────────────────────────────────
\section{Choix du frontend}
% ─────────────────────────────────────────────────────────────

Cette section présente les besoins spécifiques du frontend, évalue les
alternatives disponibles, puis justifie le choix retenu.

\subsection{Besoins spécifiques}

L'interface du projet doit couvrir deux espaces bien
distincts : d'un côté, un espace dédié au service RI
pour piloter les campagnes, de l'autre, un espace
étudiant pour consulter les destinations, suivre
sa candidature et interagir avec le chatbot.
Certaines pages affichent des données sensibles
(résultats d'affectation, GPA) et doivent être
construites côté serveur pour éviter qu'elles
ne soient accessibles avant que le rôle de
l'utilisateur soit vérifié. D'autres, comme
le suivi de candidature ou le chatbot,
nécessitent une navigation fluide et réactive.
Gérer ces deux types de pages dans le même projet,
avec une séparation fiable des accès selon le rôle,
est le principal critère de choix.

\subsection{Alternatives évaluées}

Trois frameworks ont été évalués au regard des besoins identifiés.

\subsubsection{Applications monopage — React + Vite}

React  permet de construire des interfaces
entièrement rendues dans le navigateur ~\cite{react_dev}.

\textit{Avantages :}
\begin{itemize}
  \item Simple à configurer et à démarrer
  \item Très bien documenté, avec de nombreuses
        ressources disponibles en ligne
\end{itemize}

\textit{Limites :}
\begin{itemize}
  \item Tout est chargé et affiché dans le navigateur :
        les données sensibles comme les affectations
        ou les GPA y arrivent avant même que
        le rôle de l'utilisateur soit contrôlé,
        ce qui fragilise la séparation entre
        l'espace étudiant et l'espace RI
  \item La protection des pages selon le rôle
        doit être codée manuellement, sans
        mécanisme intégré pour l'imposer
        au niveau du routing
  \item Les appels à l'API Django passent tous
        par le navigateur, ce qui expose
        les tokens d'authentification au client
\end{itemize}

\subsubsection{Angular}

Angular est un framework TypeScript développé
par Google, utilisé principalement pour des
applications d'entreprise avec beaucoup de
logique interne ~\cite{angular_dev}.

\textit{Avantages :}
\begin{itemize}
  \item Architecture très structurée, adaptée aux
        applications avec de nombreux écrans
        et règles métier
  \item Typage TypeScript natif, ce qui réduit
        les erreurs au moment du développement
  \item Mécanisme de protection des routes intégré,
        utile pour séparer les accès
        selon le rôle utilisateur
  \item Bien adapté à des interfaces
        d'administration complexes
        comme le tableau de bord RI
\end{itemize}

\textit{Limites :}
\begin{itemize}
  \item Pas de rendu côté serveur sans
        Angular Universal, qui nécessite
        une configuration supplémentaire
        non négligeable
  \item Angular demande un temps de prise en main
        important : son architecture impose
        de nombreux concepts spécifiques
        (modules, décorateurs, injection de
        dépendances) qui ralentissent
        le démarrage
  \item Moins d'exemples disponibles que
        React/Next.js pour l'intégration
        avec le SSO CAS utilisé
        par Toulouse INP
\end{itemize}

\subsubsection{Next.js 14}

Next.js est un framework basé sur React qui permet de combiner rendu côté serveur
et rendu côté client selon les besoins de chaque page ~\cite{nextjs_dev}.

\textit{Avantages :}
\begin{itemize}
  \item Rendu hybride : pages sensibles rendues côté serveur, composants
        interactifs côté client, les deux coexistant dans le même projet
  \item Middleware intégré pour bloquer l'accès aux pages selon le rôle
        utilisateur avant même leur chargement
  \item Appels à l'API Django réalisables côté serveur, sans exposer
        les tokens dans le navigateur
  \item App Router et React Server Components : chaque composant est rendu
        côté serveur par défaut, avec des layouts partagés par espace
        (RI, étudiant) qui structurent naturellement la séparation des interfaces
  \item TypeScript, Tailwind CSS et intégration SSO CAS/OIDC bien documentés
\end{itemize}

\textit{Limites :}
\begin{itemize}
  \item Configuration initiale un peu plus complexe qu'une application
        monopage classique
\end{itemize}

\subsection{Synthèse comparative}

Le tableau~\ref{tab:comparaison_frontend} récapitule les critères évalués pour chacune des trois solutions, afin de faciliter la comparaison et de justifier le choix retenu.

\renewcommand{\arraystretch}{1.35}
\begin{table}[H]
\centering
\caption{Comparaison des frameworks frontend}
\label{tab:comparaison_frontend}
\footnotesize
\begin{tabular}{|p{3.8cm}|>{\centering\arraybackslash}p{2.8cm}|>{\centering\arraybackslash}p{2.4cm}|>{\centering\arraybackslash}p{3cm}|}
\hline
\rowcolor{gray!20}
\textbf{Critère} & \textbf{React + Vite} & \textbf{Angular} & \textbf{Next.js 14} \\
\hline
Rendu côté serveur (SSR) & \xmark & Via Angular Universal (config requise) & \cmark \\
\hline
Protection des routes par rôle & Manuelle & Mécanisme intégré & Middleware intégré \\
\hline
Données sensibles côté serveur & \xmark & \xmark & \cmark \\
\hline
TypeScript natif & Optionnel & \cmark & \cmark \\
\hline
Intégration SSO CAS/OIDC & \cmark & Peu d'exemples avec le SSO CAS de Toulouse INP & \cmark \\
\hline
Rendu hybride SSR + CSR & \xmark & \xmark & \cmark \\
\hline
\end{tabular}
\end{table}

\subsection{Décision}

\textbf{Next.js 14} est retenu. Son modèle de React Server
Components permet de gérer naturellement les deux espaces du projet : les pages
sensibles du tableau de bord RI sont rendues côté serveur avant toute vérification
de rôle, tandis que les composants interactifs de l'espace étudiant s'exécutent côté client. Angular pourrait, en théorie, répondre au même besoin grâce à Angular Universal. Mais avec le temps disponible pour ce projet, Next.js permettait d'arriver au même résultat plus rapidement et avec moins de configuration à mettre en place.


% ─────────────────────────────────────────────────────────────
\section{Choix du backend}
% ─────────────────────────────────────────────────────────────

Cette section présente les besoins spécifiques du backend, évalue les
alternatives disponibles, puis justifie le choix retenu.

\subsection{Besoins spécifiques}

Le backend de ce projet ne se limite pas à une simple gestion de données. Il doit orchestrer de nombreuses entités métier liées entre elles et gérer des processus parfois complexes, comme le suivi des candidatures, des affectations ou des mobilités étudiantes.

L'application doit également assurer plusieurs fonctions essentielles : gérer les droits d'accès selon les profils (étudiants et administrateurs RI), conserver une trace des actions importantes via un journal d'audit, accompagner l'évolution de la base de données pendant le développement, et s'intégrer au système d'authentification institutionnel CAS/OIDC de Toulouse INP.

Dans ce contexte, le choix du framework ne repose pas uniquement sur les performances, mais aussi sur la qualité des outils qu'il apporte pour structurer, sécuriser et faire évoluer l'application.

\subsection{Choix du framework backend}

Cette sous-section évalue les frameworks backend candidats au regard des besoins définis ci-dessus, puis justifie le choix retenu.

\subsubsection{Alternatives évaluées}

Trois frameworks ont été évalués en regard des besoins identifiés précédemment.

\paragraph{Node.js}
Environnement d'exécution JavaScript côté serveur,
reconnu pour ses performances sur les applications
qui traitent beaucoup de requêtes simultanées ~\cite{nodejs_dev}.

\textit{Avantages :}
\begin{itemize}
  \item Bonnes performances sur les traitements
        non bloquants
  \item Grand nombre de bibliothèques disponibles
  \item Gestion native des opérations asynchrones
\end{itemize}

\textit{Limites :}
\begin{itemize}
    \item Aucun outil natif pour la base de données, l'administration ou la gestion des droits : il faut tout assembler manuellement à partir de bibliothèques séparées, ce qui alourdit la configuration et la maintenance.
    \item Des bibliothèques de machine à états comme XState existent et sont matures, mais rien d'équivalent à \texttt{django-fsm} n'est directement intégré à un ORM Node.js : il faudrait câbler soi-même la protection des transitions et la traçabilité fine des modifications au niveau des modèles, ce qui est indispensable dans ce projet.
    \item L'utilisation de Node.js implique de travailler en JavaScript, alors que le reste du système est conçu en Python, ce qui introduit une complexité supplémentaire en termes de cohérence et de maintenance.
    \item La gestion des évolutions du schéma de base de données au fil des sprints demande davantage d'efforts d'organisation et de suivi, car elle n'est pas aussi intégrée nativement dans l'écosystème.
\end{itemize}

\paragraph{FastAPI}
Framework Python récent, conçu pour construire
des APIs performantes avec une validation
automatique des données ~\cite{fastapi_dev}.

\textit{Avantages :}
\begin{itemize}
  \item Très bonnes performances grâce
        au support natif des requêtes asynchrones
  \item Validation automatique des données entrantes
  \item Documentation de l'API générée automatiquement
\end{itemize}

\textit{Limites :}
\begin{itemize}
  \item Pas d'outil intégré pour la base de données,
        l'administration ou la gestion des droits
  \item Les cycles de vie des entités,
        le journal d'audit et les permissions
        seraient à construire de zéro
  \item L'outil de migration de schéma demande davantage de configuration et d'intervention manuelle, ce qui peut ralentir les cycles d'itération en développement.
  \item FastAPI est surtout adapté à des APIs légères et rapides, et moins à des applications contenant beaucoup de logique métier.
\end{itemize}

\paragraph{Django}
Framework Python complet, pensé pour développer
rapidement des applications web avec beaucoup
de logique interne ~\cite{django_dev}.

\textit{Avantages :}
\begin{itemize}
  \item Gestion de la base de données,
        authentification, permissions et interface
        d'administration fournis nativement,
        sans aucune configuration supplémentaire
  \item \texttt{django-fsm} permet de définir
        des cycles de vie sécurisés sur chaque entité,
        avec une option qui empêche toute
        modification directe du statut
        en dehors des transitions autorisées
  \item \texttt{django-auditlog} enregistre
        automatiquement toutes les modifications
        avec l'auteur et la date, ce qui est indispensable
        pour les décisions réglementaires
  \item La gestion des évolutions de schéma
        est intégrée nativement, ce qui
        facilite les changements à chaque sprint
  \item S'intègre facilement avec le SSO CAS/OIDC
        via des bibliothèques existantes et maintenues
\end{itemize}

\textit{Limites :}
\begin{itemize}
  \item Django était historiquement limité
        sur les traitements asynchrones,
        mais cette lacune a été corrigée
        depuis la version 5
  \item Légèrement moins performant que FastAPI
        sur des calculs très complexes,
        ce qui peut être contourné en passant
        par des requêtes SQL directes
        pour les agrégations de durées CTI
\end{itemize}

\subsubsection{Synthèse comparative}

Le tableau~\ref{tab:comparaison_backend} récapitule les critères évalués pour chacune des trois solutions, afin de faciliter la comparaison et de justifier le choix retenu.

\renewcommand{\arraystretch}{1.35}
\begin{table}[H]
\centering
\caption{Comparaison des frameworks backend}
\label{tab:comparaison_backend}
\footnotesize
\begin{tabular}{|p{5cm}|>{\centering\arraybackslash}p{2.4cm}|>{\centering\arraybackslash}p{2.4cm}|>{\centering\arraybackslash}p{3.2cm}|}
\hline
\rowcolor{gray!20}
\textbf{Critère} & \textbf{Node.js} & \textbf{FastAPI} & \textbf{Django 5} \\
\hline
ORM et gestion BDD intégrés & \xmark & \xmark & \cmark \\
\hline
Interface d'administration & \xmark & \xmark & \cmark \\
\hline
Machines à états (FSM) & \xmark & \xmark & \cmark \\
\hline
Migrations de schéma & Manuel & nécessite config & \cmark \\
\hline
Authentification SSO & Manuel & Manuel & \cmark \\
\hline
Journal d'audit automatique & Manuel & Manuel & \cmark \\
\hline
Support ASGI / async & \cmark & \cmark & \cmark \\
\hline
\end{tabular}
\end{table}

\subsubsection{Décision}
\textbf{Django 5} est retenu car le projet repose sur de nombreuses entités liées entre elles, avec des cycles de vie distincts, un journal d'audit obligatoire et une gestion fine des droits entre deux profils utilisateurs.

Utiliser une autre solution aurait impliqué de reconstruire une grande partie de ces mécanismes.

Django permet de disposer de ces fonctionnalités directement, ce qui est un avantage important dans ce projet.

% Pour le déploiement en production, la configuration retenue est \textbf{Gunicorn} avec l'option \texttt{-k uvicorn.workers.UvicornWorker} : Gunicorn orchestre les workers et gère les signaux système, tandis que UvicornWorker assure le support natif de l'interface ASGI requise par Django Ninja pour les endpoints asynchrones.

% ─────────────────────────────────────────────────────────────
\subsection{Choix de la couche API}

Une fois le framework backend retenu, il reste à choisir la bibliothèque
la mieux adaptée pour construire l'API. Cette section évalue les deux
solutions compatibles avec Django et justifie le choix retenu.

\subsubsection{Besoins spécifiques}

L'API doit permettre les échanges entre Next.js
et Django, valider des structures de données
parfois complexes (vœux ordonnés, quotas par
département, rapports d'import), générer
automatiquement une documentation utilisable
par le développeur frontend, et répondre
de façon non bloquante aux requêtes qui
déclenchent des traitements longs
(lancement de l'algorithme, recalcul).

\subsubsection{Alternatives évaluées}

Deux solutions ont été évaluées pour exposer l'API entre le frontend Next.js et le backend Django.

\paragraph{Django REST Framework (DRF)}
Django REST Framework est la bibliothèque la plus utilisée pour construire
des APIs avec Django. Elle fournit un ensemble complet d'outils pour
la sérialisation, la validation et la gestion des permissions~\cite{drf_dev}.

\textit{Avantages :}
\begin{itemize}
  \item Très mature et bien documenté,
        avec une grande communauté
  \item Outillage complet : gestion des tokens JWT,
        filtres, pagination, génération de
        documentation via \texttt{drf-spectacular}
  \item Fonctionne parfaitement avec l'ORM
        et les permissions Django
\end{itemize}

\textit{Limites :}
\begin{itemize}
  \item Le support des requêtes asynchrones
        est partiel, ce qui pose problème pour
        les endpoints qui déclenchent le lancement
        de l'algorithme ou les synchronisations
        avec MoveON, Pégase et Eudonet
  \item Les sérialiseurs DRF demandent beaucoup
        de code répétitif, surtout quand on
        travaille avec plusieurs entités
  \item La documentation de l'API nécessite
        une configuration externe séparée,
        à maintenir en parallèle
\end{itemize}

\paragraph{Django Ninja}
Framework API pour Django, basé sur une validation des données structurée et typée via Pydantic ~\cite{djangoninja_dev}.

\textit{Avantages :}
\begin{itemize}
    \item Toutes les routes supportent nativement les traitements asynchrones : un endpoint peut lancer un calcul ou déclencher une synchronisation sans bloquer le serveur, en déléguant l'exécution à un système de traitement en arrière-plan.
  \item La validation des données via Pydantic
        est plus concise que DRF
  \item La documentation de l'API est générée
        automatiquement, sans aucune
        configuration supplémentaire
  \item Conserve entièrement l'écosystème Django :
        ORM, administration, cycles de vie,
        permissions
\end{itemize}

\textit{Limites :}
\begin{itemize}
  \item Plus récent que DRF, avec moins
        d'extensions disponibles
  \item Moins de ressources en ligne
        pour des cas très spécifiques
\end{itemize}

\subsubsection{Synthèse comparative}

Le tableau~\ref{tab:comparaison_api} récapitule les critères évalués pour chacune des deux
solutions, afin de faciliter la comparaison et de justifier le choix retenu.

\renewcommand{\arraystretch}{1.35}
\begin{table}[H]
\centering
\caption{Comparaison DRF et Django Ninja}
\label{tab:comparaison_api}
\footnotesize
\begin{tabular}{|p{5cm}|>{\centering\arraybackslash}p{3.2cm}|>{\centering\arraybackslash}p{3.2cm}|}
\hline
\rowcolor{gray!20}
\textbf{Critère} & \textbf{DRF} & \textbf{Django Ninja} \\
\hline
Validation Pydantic (typage strict) & \xmark & \cmark \\
\hline
Réduction du code répétitif & \xmark & \cmark \\
\hline
Documentation auto-générée & \xmark & \cmark \\
\hline
Intégration ORM Django & \cmark & \cmark \\
\hline
Support async natif (vues) & Non natif, configuration requise & \cmark \\
\hline
Maturité et communauté & \cmark & Framework récent, moins de ressources \\
\hline
\end{tabular}
\end{table}

\subsubsection{Décision}
\textbf{Django Ninja} est retenu. L'argument principal est la validation
par schémas Pydantic : les entrées et sorties de chaque endpoint sont
typées explicitement, ce qui réduit nettement le code répétitif de
sérialisation par rapport à DRF, en particulier sur les nombreuses
entités du domaine (accords, quotas, vœux, affectations). La
documentation OpenAPI générée automatiquement à partir de ces schémas
est également appréciable pour un projet développé avec un frontend
séparé.

Le support natif de l'asynchrone n'est pas le critère décisif : les
endpoints critiques du projet (déclenchement de l'algorithme, recalcul
des quotas, synchronisations ETL) répondent immédiatement en
déléguant le traitement à une tâche Django-Q2, ce qu'une vue
synchrone classique fait tout aussi bien puisque c'est la mise en
file d'attente qui est asynchrone, pas la vue elle-même. Ce support
reste néanmoins un atout pour d'éventuels futurs endpoints qui
appelleraient directement une API externe (Eudonet, MoveON, Pégase)
sans passer par une tâche de fond.

\medskip
\noindent La stack backend retenue est donc :
\textbf{Django 5 + Django Ninja}.

% ─────────────────────────────────────────────────────────────
\section{Choix de la base de données}
% ─────────────────────────────────────────────────────────────

Cette section présente les besoins spécifiques liés au stockage des données,
évalue les alternatives disponibles, puis justifie le choix retenu.

\subsection{Besoins spécifiques}

La base de données doit supporter un modèle avec plusieurs entités reliées
entre elles par de nombreuses clés étrangères, stocker les réponses brutes
des APIs externes avant transformation (\texttt{RawImport}), et garantir
la cohérence des données lors des opérations qui touchent plusieurs entités
à la fois (transitions de cycle de vie, affectations, journal d'audit).

\subsection{Alternatives évaluées}

Trois systèmes de gestion de base de données ont été évalués en regard
des besoins identifiés précédemment.

\subsubsection{MySQL}

MySQL est un système de gestion de base de données relationnelle très répandu,
bien supporté par l'écosystème Django et largement utilisé en
production ~\cite{mysql_dev}.

\textit{Avantages :}
\begin{itemize}
  \item Système très répandu, avec une grande communauté et une documentation abondante
  \item Bonne intégration avec l'écosystème Django
  \item Largement utilisé en production, avec des outils de sauvegarde matures
\end{itemize}

\textit{Limites :}
\begin{itemize}
  \item Support moins adapté pour le stockage et l'exploitation de données
        au format JSON, notamment pour les imports bruts (\texttt{RawImport})
  \item Fonctionnalités limitées pour gérer des mécanismes de file d'attente
        de tâches
  \item Moins adapté aux opérations transactionnelles complexes impliquant
        plusieurs entités liées
\end{itemize}

\subsubsection{MongoDB}

MongoDB est une base de données orientée documents, conçue pour offrir
une grande flexibilité dans la structure des données et particulièrement
adaptée au stockage JSON ~\cite{mongodb_dev}.

\textit{Avantages :}
\begin{itemize}
  \item Grande flexibilité dans la structure des données
  \item Adapté au stockage de documents JSON
  \item Bonnes performances sur les lectures de documents non structurés
\end{itemize}

\textit{Limites :}
\begin{itemize}
  \item Le modèle de données du projet repose fortement sur des relations
        entre entités (quotas, affectations, mobilités), ce qui rend
        l'absence de relations structurées moins adaptée
  \item Intégration limitée avec les outils ORM utilisés dans le projet
  \item Absence de support direct pour certains mécanismes de traitement
        asynchrone utilisés dans l'architecture
\end{itemize}

\subsubsection{PostgreSQL 16}

PostgreSQL est un système de gestion de base de données relationnelle
open source, reconnu pour sa robustesse, la richesse de ses fonctionnalités
et son respect strict des standards SQL ~\cite{postgresql_dev}.

\textit{Avantages :}
\begin{itemize}
  \item Gestion robuste des transactions, adaptée aux opérations impliquant
        plusieurs entités liées
  \item Support natif du format JSON avec indexation, utile pour le stockage
        des données brutes (\texttt{RawImport})
  \item Possibilité d'utiliser la base comme support de file d'attente
        pour les traitements asynchrones
  \item Bon niveau d'intégration avec les outils de développement du projet
  \item Outils de sauvegarde et de restauration fiables
\end{itemize}

\textit{Limites :}
\begin{itemize}
  \item Mise en place initiale plus complète, nécessitant une configuration
        structurée dès le début du projet
\end{itemize}

\subsection{Synthèse comparative}

Le tableau~\ref{tab:comparaison_bdd} récapitule les critères évalués pour chacune des trois
solutions, afin de faciliter la comparaison et de justifier le choix retenu.

\renewcommand{\arraystretch}{1.35}
\begin{table}[H]
\centering
\caption{Comparaison des systèmes de gestion de base de données}
\label{tab:comparaison_bdd}
\footnotesize
\begin{tabular}{|p{4cm}|>{\centering\arraybackslash}p{2.6cm}|>{\centering\arraybackslash}p{2.4cm}|>{\centering\arraybackslash}p{2.8cm}|}
\hline
\rowcolor{gray!20}
\textbf{Critère} & \textbf{MySQL / MariaDB} & \textbf{MongoDB} & \textbf{PostgreSQL 16} \\
\hline
Transactions ACID multi-tables & \cmark & ACID ajouté en v4.0, moins fiable qu'un SGBD relationnel & \cmark \\
\hline
Support JSONB natif + indexation & JSON simple, pas d'indexation avancée & \cmark & \cmark \\
\hline
File d'attente de tâches & \xmark & \xmark & \cmark \\
\hline
Intégration ORM Django & \cmark & Support via PyMongo sans migration Django & \cmark \\
\hline
Extension pgcrypto (chiffrement) & \xmark & \xmark & \cmark \\
\hline
Outils de sauvegarde fiables & \cmark & \cmark & \cmark \\
\hline
\end{tabular}
\end{table}

\subsection{Décision}

\textbf{PostgreSQL 16} est retenu. Il répond de manière cohérente à l'ensemble
des besoins identifiés, à savoir la gestion des relations, le support du JSON natif,
les traitements asynchrones et le chiffrement, dans un seul service, sans introduire
de composant supplémentaire dans l'infrastructure.

% ─────────────────────────────────────────────────────────────
\section{Choix du format de stockage des imports bruts}
% ─────────────────────────────────────────────────────────────

Cette section précise comment sont stockées les réponses brutes des APIs
externes avant leur traitement, en évaluant les alternatives disponibles
et en justifiant le format retenu.

\subsection{Besoins spécifiques}

L'entité \texttt{RawImport} doit conserver la réponse complète de chaque
appel d'API externe (MoveON, Pégase, Eudonet) avant toute transformation.
Ces données sont semi-structurées et variables d'une source à l'autre.
Le format retenu doit permettre de relire, d'interroger et de tracer ces
imports tout en restant cohérent avec le reste du modèle de données.

\subsection{Alternatives évaluées}

\subsubsection{Fichiers plats (JSON sur le système de fichiers)}

Un fichier plat est un fichier texte brut qui stocke des données sans
structure hiérarchique ni mécanisme d'indexation. Dans ce contexte, chaque
réponse d'API serait sérialisée en JSON et enregistrée dans un répertoire
dédié sur le serveur, un fichier par import.

\textit{Avantages :}
\begin{itemize}
  \item Mise en place immédiate, sans configuration particulière
  \item Lecture directe par n'importe quel éditeur de texte
\end{itemize}

\textit{Limites :}
\begin{itemize}
  \item Aucune possibilité d'interroger le contenu avec des requêtes SQL
  \item Intégrité référentielle avec les autres entités non garantie
  \item Gestion des accès concurrents et des transactions non assurée
  \item Sauvegarde dissociée du reste de la base de données
\end{itemize}

\subsubsection{Base NoSQL dédiée (MongoDB)}

Les imports bruts sont stockés dans une base documentaire séparée,
spécialisée dans le stockage de documents JSON.

\textit{Avantages :}
\begin{itemize}
  \item Très bien adapté au stockage de documents JSON hétérogènes
  \item Bonne performance sur les lectures de documents non structurés
\end{itemize}

\textit{Limites :}
\begin{itemize}
  \item Service supplémentaire à déployer, superviser et sauvegarder
  \item Aucune relation possible avec les entités de la base relationnelle
  \item Surcoût d'infrastructure injustifié pour ce seul cas d'usage
\end{itemize}

\subsubsection{Type JSONB dans PostgreSQL}

Les réponses brutes sont stockées dans une colonne de type \textbf{JSONB}
au sein de la table \texttt{RawImport}, dans la même base PostgreSQL.

\textit{Avantages :}
\begin{itemize}
  \item Stockage dans la base existante, sans service supplémentaire
  \item Indexation du contenu JSON pour des requêtes efficaces
  \item Jointures possibles avec les autres tables du modèle
  \item Mêmes garanties ACID et même périmètre de sauvegarde que le reste
\end{itemize}

\textit{Limites :}
\begin{itemize}
  \item Moins flexible que MongoDB pour des structures JSON très hétérogènes,
        ce qui reste sans impact ici compte tenu du faible nombre de sources
\end{itemize}

\subsection{Synthèse comparative}

\renewcommand{\arraystretch}{1.35}
\begin{table}[H]
\centering
\caption{Comparaison des formats de stockage pour \texttt{RawImport}}
\label{tab:comparaison_rawimport}
\footnotesize
\begin{tabular}{|p{4.5cm}|>{\centering\arraybackslash}p{2.5cm}|>{\centering\arraybackslash}p{2.5cm}|>{\centering\arraybackslash}p{2.8cm}|}
\hline
\rowcolor{gray!20}
\textbf{Critère} & \textbf{Fichiers plats} & \textbf{MongoDB} & \textbf{JSONB (PostgreSQL)} \\
\hline
Requêtes SQL sur le contenu & \xmark & \xmark & \cmark \\
\hline
Intégrité référentielle & \xmark & \xmark & \cmark \\
\hline
Garanties ACID & \xmark & Transactions ACID ajoutées en v4.0, moins fiables que PostgreSQL & \cmark \\
\hline
Aucun service supplémentaire & \cmark & \xmark & \cmark \\
\hline
Indexation du contenu JSON & \xmark & \cmark & \cmark \\
\hline
Sauvegarde unifiée & \xmark & \xmark & \cmark \\
\hline
\end{tabular}
\end{table}

\subsection{Décision}

Le type \textbf{JSONB} de PostgreSQL est retenu pour le stockage des imports
bruts. Il évite toute fragmentation de l'infrastructure en réutilisant la base
déjà en place, tout en offrant les capacités d'interrogation nécessaires sur
les données semi-structurées. Les fichiers plats ne permettent pas de garantir
la cohérence avec le reste du modèle, et l'ajout de MongoDB n'apporterait
aucun avantage fonctionnel pour ce seul cas d'usage.

% ─────────────────────────────────────────────────────────────
\section{Choix du gestionnaire de tâches asynchrones}
% ─────────────────────────────────────────────────────────────

Cette section présente les besoins en traitement asynchrone du projet,
évalue les solutions disponibles, puis justifie le choix retenu.

\subsection{Besoins spécifiques}

Certaines opérations du système ne peuvent pas s'exécuter de façon
synchrone sans bloquer l'interface utilisateur. Deux catégories de
tâches ont été identifiées.

La première catégorie regroupe les tâches déclenchées par des événements :
le lancement de l'algorithme d'affectation, le recalcul des durées
après modification d'un état, et le traitement des imports qui peuvent
durer plusieurs secondes.

La seconde catégorie regroupe les tâches planifiées : les synchronisations
nocturnes avec MoveON, Pégase et Eudonet, ainsi que la suppression
automatique des imports bruts anciens.

Le volume reste modéré : quelques dizaines de tâches par jour en
fonctionnement normal, quelques centaines lors des périodes de traitement,
sans besoin de traitement intensif en continu.

\subsection{Alternatives évaluées}

Deux solutions ont été évaluées pour la gestion des tâches asynchrones
dans l'écosystème Django.

\subsubsection{Celery + Redis}

Celery est la bibliothèque Python de référence pour l'exécution de tâches
en arrière-plan. Elle s'appuie sur un broker de messages externe,
généralement Redis ou RabbitMQ, pour acheminer les tâches vers les
workers ~\cite{celery_dev} ~\cite{redis_dev}.

\textit{Avantages :}
\begin{itemize}
  \item Très mature, bien documenté, avec une grande communauté
  \item Capable de traiter un très grand volume de tâches en parallèle
  \item Compatible avec plusieurs brokers (Redis, RabbitMQ)
  \item Outil de supervision dédié (Flower)
\end{itemize}

\textit{Limites :}
\begin{itemize}
  \item Trois services supplémentaires à déployer et maintenir :
        le worker Celery, le planificateur Celery Beat et Redis
  \item Redis consomme entre 50 et 200~Mo de mémoire pour un rôle
        que PostgreSQL, déjà présent, peut assurer seul
  \item Complexité disproportionnée par rapport au volume réel du projet
\end{itemize}

\subsubsection{Django-Q2 + PostgreSQL}

Django-Q2 est une bibliothèque de gestion de tâches asynchrones conçue
spécifiquement pour Django. Elle utilise directement PostgreSQL comme
file d'attente, sans nécessiter de broker externe ~\cite{djangoq2_dev}.

\textit{Avantages :}
\begin{itemize}
  \item Aucun service supplémentaire : PostgreSQL joue le rôle
        de file d'attente sans configuration additionnelle
  \item Planificateur de tâches intégré, adapté aux synchronisations
        nocturnes et aux suppressions programmées
  \item Supervision accessible depuis l'interface d'administration
        Django existante, sans outil tiers
  \item Seulement deux services à opérer au lieu de cinq avec Celery
\end{itemize}

\textit{Limites :}
\begin{itemize}
  \item Moins performant que Celery pour des volumes très élevés,
        sans incidence pour ce projet
  \item Communauté plus restreinte que Celery
\end{itemize}

\subsection{Synthèse comparative}

Le tableau~\ref{tab:comparaison_async} récapitule les critères évalués pour chacune des deux
solutions, afin de faciliter la comparaison et de justifier le choix retenu.

\renewcommand{\arraystretch}{1.35}
\begin{table}[H]
\centering
\caption{Comparaison des gestionnaires de tâches asynchrones}
\label{tab:comparaison_async}
\footnotesize
\begin{tabular}{|p{5.5cm}|>{\centering\arraybackslash}p{3cm}|>{\centering\arraybackslash}p{3cm}|}
\hline
\rowcolor{gray!20}
\textbf{Critère} & \textbf{Celery + Redis} & \textbf{Django-Q2 + PostgreSQL} \\
\hline
Aucun service supplémentaire & \xmark & \cmark \\
\hline
Planificateur intégré & Celery Beat disponible mais service séparé & \cmark \\
\hline
Supervision via admin Django & \xmark & \cmark \\
\hline
Adapté au volume du projet & Surdimensionné pour ce volume de tâches & \cmark \\
\hline
Maturité et communauté & \cmark & Communauté plus restreinte que Celery \\
\hline
Intégration native avec Django & Nécessite une configuration supplémentaire & \cmark \\
\hline
\end{tabular}
\end{table}

\subsection{Décision}

\textbf{Django-Q2 avec PostgreSQL} est retenu. Il réutilise la base de
données déjà en place comme file d'attente, ce qui évite d'introduire
un service supplémentaire dans l'infrastructure. Le volume de tâches
du projet ne justifie pas la complexité opérationnelle de Celery. La
supervision des traitements depuis l'interface d'administration Django
existante simplifie également le suivi quotidien pour le service RI.


% ─────────────────────────────────────────────────────────────
\section{Choix du stockage de fichiers}
% ─────────────────────────────────────────────────────────────

Cette section présente les besoins en stockage de fichiers du projet,
évalue les solutions disponibles, puis justifie le choix retenu.

\subsection{Besoins spécifiques}

Le projet requiert un mécanisme de dépôt et de stockage de fichiers pour
les justificatifs soumis par les étudiants lors de la déclaration de
mobilités complémentaires. Ces fichiers doivent être stockés de manière
sécurisée et chiffrée, accessibles uniquement par le service RI et
l'étudiant concerné, avec une suppression automatique conforme au RGPD.

\subsection{Alternatives évaluées}

Deux approches ont été évaluées pour répondre à ces besoins.

\subsubsection{Système de fichiers local}

Le stockage local consiste à enregistrer les fichiers directement sur
le système de fichiers du serveur hébergeant l'application Django,
dans un répertoire dédié.

\textit{Avantages :}
\begin{itemize}
  \item Simple à configurer, sans dépendance supplémentaire
\end{itemize}

\textit{Limites :}
\begin{itemize}
  \item Aucun chiffrement natif des fichiers stockés
  \item Non portable entre environnements (développement, staging,
        production)
  \item Difficile à sauvegarder et à migrer indépendamment
        de la base de données
\end{itemize}

\subsubsection{MinIO (stockage objet S3-compatible)}

MinIO est une solution de stockage objet open source, auto-hébergeable
et compatible avec l'API S3 d'Amazon ~\cite{minio_dev}. Elle s'intègre
nativement dans une infrastructure Docker et ne nécessite aucun service
cloud externe.

\textit{Avantages :}
\begin{itemize}
  \item Stockage objet auto-hébergé, compatible avec l'API S3 standard
  \item Chiffrement côté serveur (SSE-S3) des fichiers stockés
  \item Génération d'URLs pré-signées pour un accès temporaire
        et sécurisé, sans exposer les fichiers directement
  \item Déployable comme container Docker interne,
        sans aucune dépendance externe ni service cloud
  \item Politique de rétention et suppression automatique configurables,
        conformes aux exigences RGPD
\end{itemize}

\textit{Limites :}
\begin{itemize}
  \item Configuration initiale plus complexe qu'un stockage local
\end{itemize}

\subsection{Synthèse comparative}

Le tableau~\ref{tab:comparaison_stockage} récapitule les critères évalués pour chacune des deux
solutions, afin de faciliter la comparaison et de justifier le choix retenu.

\renewcommand{\arraystretch}{1.35}
\begin{table}[H]
\centering
\caption{Comparaison des solutions de stockage de fichiers}
\label{tab:comparaison_stockage}
\footnotesize
\begin{tabular}{|p{5.5cm}|>{\centering\arraybackslash}p{3cm}|>{\centering\arraybackslash}p{3cm}|}
\hline
\rowcolor{gray!20}
\textbf{Critère} & \textbf{Fichiers local} & \textbf{MinIO} \\
\hline
Chiffrement natif (SSE-S3) & \xmark & \cmark \\
\hline
URLs pré-signées (accès sécurisé) & \xmark & \cmark \\
\hline
Portabilité multi-environnements & \xmark & \cmark \\
\hline
Suppression automatique (RGPD) & \xmark & \cmark \\
\hline
Déploiement Docker intégré & \xmark & \cmark \\
\hline
Simplicité de configuration & \cmark & Plus complexe qu'un stockage local \\
\hline
\end{tabular}
\end{table}

\subsection{Décision}

\textbf{MinIO} est retenu. MinIO peut être déployé selon deux modes :
\textbf{standalone} (instance unique sur un seul serveur) ou
\textbf{distributed} (plusieurs nœuds pour la haute disponibilité et
la tolérance aux pannes). Le mode distribué est destiné aux
infrastructures multi-serveurs nécessitant une réplication des données.

Dans le cadre de ce projet, le déploiement cible un \textbf{serveur
unique on-premise} au sein de l'établissement. Le mode standalone
est donc retenu : il correspond exactement à cette topologie, sans
surcharger l'infrastructure avec des mécanismes de distribution
inutiles. MinIO répond ainsi directement aux exigences de sécurité
(chiffrement SSE-S3) et de conformité RGPD (rétention et suppression
automatiques), tout en s'intégrant de façon transparente dans
Docker Compose sans aucun service externe.


% ─────────────────────────────────────────────────────────────
\section{Couche de routage des requêtes}
% ─────────────────────────────────────────────────────────────

Cette section présente les besoins en matière de routage réseau,
évalue les deux approches disponibles, puis justifie le choix retenu.

\subsection{Besoins spécifiques}

Le système doit disposer d'un composant chargé de recevoir les
requêtes HTTP entrantes, d'assurer la terminaison SSL, de router
le trafic vers le frontend Next.js et l'API Django Ninja, et de
servir directement les fichiers statiques sans les faire transiter par l'application. Ce composant
sera administré par
les équipes IT de Toulouse INP ; sa configuration doit donc être
lisible et maintenable sans expertise Docker particulière.

\subsection{Alternatives évaluées}

Deux approches ont été évaluées pour assurer ce rôle.

\subsubsection{Reverse proxy}

Un reverse proxy est un composant qui se place entre les clients et
les services internes. Il reçoit les requêtes HTTP entrantes, assure
la terminaison SSL, les achemine vers le service approprié selon des
règles de routage, et sert les fichiers statiques directement sans
solliciter l'application.

\textit{Avantages :}
\begin{itemize}
  \item Composant léger, dédié au routage et à la terminaison SSL
  \item Configuration statique (Nginx) ou dynamique (Traefik) selon
        le besoin
  \item Bien adapté à une architecture monolithique sur un seul serveur
  \item Maturité et documentation abondantes pour les deux outils
\end{itemize}

\textit{Limites :}
\begin{itemize}
  \item Ne centralise pas l'authentification ni le contrôle de débit
        (ce qui n'est pas requis ici)
\end{itemize}

\subsubsection{API Gateway}

Une API Gateway est un composant qui, en plus du routage, centralise
la gestion de l'authentification, du contrôle de débit et de la
transformation des requêtes. Elle est conçue pour exposer plusieurs
microservices à des clients externes variés.

\textit{Avantages :}
\begin{itemize}
  \item Gestion centralisée de l'authentification et du contrôle de débit
  \item Adapté aux architectures microservices exposées à l'externe
\end{itemize}

\textit{Limites :}
\begin{itemize}
  \item Complexité et surcoût injustifiés pour une architecture
        déployée sur un seul serveur
  \item Le projet consomme des APIs externes, il n'en expose pas
        à des clients variés ; ce cas d'usage ne s'applique pas
\end{itemize}

\subsection{Synthèse comparative}

Le tableau~\ref{tab:comparaison_proxy} récapitule les critères évalués pour chacune des
deux approches, afin de faciliter la comparaison et de justifier
le choix retenu.

\renewcommand{\arraystretch}{1.35}
\begin{table}[H]
\centering
\caption{Comparaison des approches de routage des requêtes}
\label{tab:comparaison_proxy}
\footnotesize
\begin{tabular}{|p{5.5cm}|>{\centering\arraybackslash}p{3cm}|>{\centering\arraybackslash}p{3cm}|}
\hline
\rowcolor{gray!20}
\textbf{Critère} & \textbf{Reverse proxy} & \textbf{API Gateway} \\
\hline
Adapté à un seul serveur & \cmark & \xmark \\
\hline
Terminaison SSL & \cmark & \cmark \\
\hline
Routage vers services internes & \cmark & \cmark \\
\hline
Gestion auth. centralisée & \xmark & \cmark \\
\hline
Complexité d'exploitation & Faible & Élevée \\
\hline
Pertinent pour ce projet & \cmark & \xmark \\
\hline
\end{tabular}
\end{table}

\subsection{Décision}

Le \textbf{reverse proxy} est retenu. L'architecture est déployée sur
un seul serveur et le système consomme des APIs externes sans en
exposer à des clients variés. Une API Gateway introduirait une
complexité sans aucun bénéfice fonctionnel dans ce contexte.

\subsection{Choix de l'implémentation : Nginx vs Traefik}

Une fois le reverse proxy retenu, il reste à choisir l'outil
d'implémentation. Cette partie évalue les deux solutions
les plus utilisées dans l'écosystème Docker et justifie le choix retenu.

\subsubsection{Besoins spécifiques}

L'outil retenu doit assurer la terminaison SSL, le routage vers
le frontend Next.js et l'API Django Ninja, et la distribution des
fichiers statiques. Sa configuration doit être lisible et auditable
par les équipes IT de Toulouse INP qui administreront le serveur,
sans expertise Docker avancée.

\subsubsection{Alternatives évaluées}

Deux outils open source ont été comparés pour implémenter le reverse proxy.

\paragraph{Nginx}

Nginx est un serveur web et reverse proxy très répandu, reconnu pour
sa stabilité et la lisibilité de ses fichiers de
configuration~\cite{nginx_dev}.

\textit{Avantages :}
\begin{itemize}
  \item Très stable et bien documenté, avec un comportement prévisible
  \item Configuration dans des fichiers texte simples, lisibles
        et faciles à auditer par n'importe quel administrateur système
  \item Gestion SSL/TLS fiable, excellentes performances sur
        le routage des requêtes
\end{itemize}

\textit{Limites :}
\begin{itemize}
  \item La gestion des certificats SSL nécessite un outil externe
        (Certbot)
  \item Configuration manuelle à mettre à jour lors de changements
        de services
\end{itemize}

\paragraph{Traefik}

Traefik est un reverse proxy conçu pour les environnements
conteneurisés. Il se configure dynamiquement à partir des labels
des conteneurs Docker actifs~\cite{traefik_dev}.

\textit{Avantages :}
\begin{itemize}
  \item Configuration automatique à partir des conteneurs Docker actifs
  \item Gestion native des certificats SSL via Let's Encrypt
  \item Interface de supervision intégrée
\end{itemize}

\textit{Limites :}
\begin{itemize}
  \item Configuration déclarative moins lisible pour des équipes IT
        non familières avec Docker
  \item Comportement moins prévisible dans un environnement
        institutionnel avec des contraintes réseau spécifiques
\end{itemize}

\subsubsection{Synthèse comparative}

Le tableau~\ref{tab:comparaison_proxy_outil} récapitule les critères évalués pour chacun des
deux outils, afin de faciliter la comparaison et de justifier le
choix retenu.

\renewcommand{\arraystretch}{1.35}
\begin{table}[H]
\centering
\caption{Comparaison des outils de reverse proxy}
\label{tab:comparaison_proxy_outil}
\footnotesize
\begin{tabular}{|p{5.5cm}|>{\centering\arraybackslash}p{3cm}|>{\centering\arraybackslash}p{3cm}|}
\hline
\rowcolor{gray!20}
\textbf{Critère} & \textbf{Nginx} & \textbf{Traefik} \\
\hline
Configuration lisible (fichiers texte) & \cmark & Labels Docker, moins lisible pour les équipes IT \\
\hline
Gestion SSL native & Nécessite Certbot pour les certificats & \cmark \\
\hline
Automatisation Docker & Configuration manuelle à chaque changement & \cmark \\
\hline
Prévisibilité en environnement institutionnel & \cmark & Moins prévisible avec des contraintes réseau institutionnelles \\
\hline
Maintenabilité par équipes IT générales & \cmark & \xmark \\
\hline
Maturité et documentation & \cmark & \cmark \\
\hline
\end{tabular}
\end{table}

\subsubsection{Décision}

\textbf{Nginx} est retenu. Le serveur sera administré par les équipes
IT de Toulouse INP. Disposer d'une configuration statique dans des
fichiers texte qu'un administrateur système peut lire et modifier
sans expertise Docker est, dans ce contexte institutionnel, plus
important que l'automatisation dynamique proposée par Traefik.

% ─────────────────────────────────────────────────────────────
\section{Synthèse des choix technologiques}
% ─────────────────────────────────────────────────────────────

\renewcommand{\arraystretch}{1.4}

\begin{longtable}{|p{3cm}|p{3.5cm}|p{8.5cm}|}
\caption{Synthèse des choix technologiques} \\
\hline
\rowcolor{gray!20}
\textbf{Composant} & \textbf{Technologie retenue} & \textbf{Justification principale} \\
\hline
\endfirsthead

\hline
\rowcolor{gray!20}
\textbf{Composant} & \textbf{Technologie retenue} & \textbf{Justification principale} \\
\hline
\endhead

Méthodologie de développement & Scrum & Développement itératif et adaptation progressive aux besoins du service RI \\
\hline

Frontend & Next.js 14 & Rendu hybride SSR/CSR, séparation sécurisée des espaces RI et étudiant \\
\hline

Backend framework & Django 5 & ORM, administration, gestion des droits et logique métier intégrés \\
\hline

Couche API & Django Ninja & Support natif de l'asynchrone, validation Pydantic, documentation automatique \\
\hline

Base de données & PostgreSQL 16 & Transactions robustes, support JSON et bonne intégration avec Django \\
\hline

Gestionnaire de tâches asynchrones & Django-Q2 + PostgreSQL & Aucun service supplémentaire et supervision intégrée à Django \\
\hline

Stockage fichiers & MinIO (S3-compatible) & Stockage objet auto-hébergé, chiffrement SSE-S3, conforme RGPD, intégré comme container Docker sans service externe \\
\hline

Reverse proxy & Nginx & Configuration simple, stable et adaptée à un environnement institutionnel \\
\hline

Conteneurisation & Docker + Docker Compose & Déploiement reproductible et simplification de l'environnement \\
\hline

Gestion des états métier & django-fsm & Transitions sécurisées et intégration directe avec Django ORM \\
\hline

Journalisation et audit & django-auditlog & Historique automatique des modifications et traçabilité \\
\hline

Authentification & CAS/OIDC & Exigence institutionnelle : SSO Toulouse INP avec détection automatique du rôle depuis le token \\
\hline

\end{longtable}


% ─────────────────────────────────────────────────────────────
\section{Stratégie de performance base de données}
% ─────────────────────────────────────────────────────────────

Cette section présente les enjeux de performance liés à la croissance
des données, évalue les deux approches disponibles, puis justifie
la stratégie retenue.

\subsection{Contexte et enjeux}

Plusieurs tables du projet, telles que les mobilités, les affectations ou les imports,
grossissent d'une campagne à l'autre. Le système doit pouvoir les
filtrer rapidement par année universitaire, que ce soit pour afficher
le tableau de bord RI, calculer les durées CTI ou lancer l'algorithme
d'affectation. Deux approches ont été évaluées pour répondre à cet
enjeu.

\subsection{Alternatives évaluées}

\subsubsection{Partitionnement par année}

Le partitionnement consiste à découper physiquement une table en
sous-tables, une par valeur de la colonne de partition, ici par
année universitaire. PostgreSQL ne lit alors que la partition
concernée lors d'une requête filtrée.

\textit{Avantages :}
\begin{itemize}
  \item Réduction significative du volume scanné pour les requêtes
        filtrées par année
  \item Efficace lorsque les tables atteignent plusieurs millions
        de lignes
\end{itemize}

\textit{Limites :}
\begin{itemize}
  \item La colonne de partitionnement doit être incluse dans toutes
        les contraintes d'unicité, ce qui peut créer des conflits
        avec les migrations Django au fil des sprints
  \item Complexité supplémentaire du schéma pour un volume qui
        ne le justifie pas
\end{itemize}

\subsubsection{Indexation ciblée}

L'indexation ciblée consiste à créer des index sur les colonnes
les plus utilisées dans les filtres, sans modifier la structure
physique des tables.

\textit{Avantages :}
\begin{itemize}
  \item Aucune modification du schéma ni des migrations Django
  \item Index déclarés directement dans les modèles, visibles
        dans le code et versionnés avec lui
  \item Transparent pour l'ORM, sans configuration PostgreSQL séparée
\end{itemize}

\textit{Limites :}
\begin{itemize}
  \item Moins efficace que le partitionnement pour des volumes
        de plusieurs millions de lignes, sans incidence ici
\end{itemize}

\subsection{Synthèse comparative}

Le tableau~\ref{tab:comparaison_perf_bdd} récapitule les critères évalués pour chacune des
deux approches, afin de faciliter la comparaison et de justifier
la stratégie retenue.

\renewcommand{\arraystretch}{1.35}
\begin{table}[H]
\centering
\caption{Comparaison des stratégies de performance base de données}
\label{tab:comparaison_perf_bdd}
\footnotesize
\begin{tabular}{|p{5.5cm}|>{\centering\arraybackslash}p{3cm}|>{\centering\arraybackslash}p{3cm}|}
\hline
\rowcolor{gray!20}
\textbf{Critère} & \textbf{Partitionnement} & \textbf{Indexation ciblée} \\
\hline
Adapté au volume du projet & \xmark & \cmark \\
\hline
Compatible avec les migrations Django & Contraintes d'unicité incompatibles avec certaines migrations & \cmark \\
\hline
Déclaration dans les modèles & \xmark & \cmark \\
\hline
Aucune configuration PostgreSQL séparée & \xmark & \cmark \\
\hline
Efficacité à très grande échelle & \cmark & Moins efficace pour plusieurs millions de lignes \\
\hline
\end{tabular}
\end{table}

\subsection{Décision}

\textbf{L'indexation ciblée} est retenue. L'ENSEEIHT traite quelques
centaines d'étudiants en mobilité par an. Même en projetant ce volume
sur dix ans, les tables restent très loin du seuil où le partitionnement
apporterait un gain mesurable. Des index sont créés sur
\texttt{academic\_year\_id} et sur la combinaison
\texttt{(academic\_year\_id, student\_id)} pour les tables les plus
sollicitées, directement dans les modèles Django concernés.

% ─────────────────────────────────────────────────────────────
\section{Choix de l'algorithme d'affectation}
% ─────────────────────────────────────────────────────────────

Cette section présente les contraintes auxquelles l'algorithme doit
répondre, évalue les approches disponibles, puis justifie le choix retenu.

\subsection{Besoins spécifiques}

L'algorithme doit attribuer à chaque étudiant une destination parmi ses
vœux, en respectant un ensemble de contraintes institutionnelles. Les
résultats ont un impact réel sur les parcours : le service RI doit pouvoir
justifier chaque affectation. L'équité entre candidats et la
reproductibilité des résultats sont des exigences essentielles.

Le classement des étudiants repose principalement sur le GPA, complété
par des règles institutionnelles. Les affectations doivent respecter
les quotas par accord de mobilité et par département, les contraintes
de financement, ainsi que la situation des étudiants étrangers dans
le classement global. Chaque étudiant exprime jusqu'à trois vœux
ordonnés ; l'algorithme traite les vœux dans l'ordre et recherche
une affectation alternative en cas d'indisponibilité.

\subsection{Alternatives évaluées}

Deux approches algorithmiques ont été évaluées pour répondre à ces besoins.

\subsubsection{Algorithme glouton}

L'algorithme glouton parcourt les candidats dans l'ordre de leur GPA
et affecte chacun à son meilleur vœu encore disponible, sans jamais
revenir sur une décision déjà prise.

\textit{Avantages :}
\begin{itemize}
  \item Simple à implémenter et à expliquer
  \item Résultat prévisible, complexité $O(n^2)$
\end{itemize}

\textit{Limites :}
\begin{itemize}
  \item Le résultat dépend de l'ordre de traitement : deux étudiants
        avec des GPA très proches peuvent se retrouver dans des
        situations très différentes selon leur position dans la liste
  \item Une légère variation de l'ordre d'entrée peut produire des
        résultats différents, ce qui est incompatible avec l'exigence
        de reproductibilité
\end{itemize}

\subsubsection{Gale-Shapley}

L'algorithme de Gale-Shapley est un algorithme
d'affectation stable qui procède par rondes successives de propositions
et de rejets différés jusqu'à convergence. Il garantit qu'aucune paire
étudiant-destination ne peut former une meilleure association que celle
retenue à l'issue du processus~\cite{GaleShapley1962}.

\textit{Avantages :}
\begin{itemize}
  \item Résultat stable : il est mathématiquement impossible qu'un
        étudiant et une destination forment une meilleure paire que
        celle retenue, quelles que soient les données
  \item Équitable : le résultat ne dépend pas de l'ordre de traitement,
        deux exécutions avec les mêmes données donnent toujours le
        même résultat
  \item Optimal pour les étudiants dans la variante Student-Proposing :
        chacun obtient la meilleure destination possible compte tenu
        de son classement et de ses vœux
  \item Complexité $O(n^2)$, tout à fait adaptée à quelques centaines
        d'étudiants
\end{itemize}

\textit{Limites :}
\begin{itemize}
  \item Légèrement plus complexe à implémenter que l'algorithme glouton
\end{itemize}

\subsection{Synthèse comparative}

Le tableau~\ref{tab:comparaison_algo} récapitule les critères évalués pour chacune des deux
approches, afin de faciliter la comparaison et de justifier le choix retenu.

\renewcommand{\arraystretch}{1.35}
\begin{table}[H]
\centering
\caption{Comparaison des algorithmes d'affectation}
\label{tab:comparaison_algo}
\footnotesize
\begin{tabular}{|p{5.5cm}|>{\centering\arraybackslash}p{3cm}|>{\centering\arraybackslash}p{3cm}|}
\hline
\rowcolor{gray!20}
\textbf{Critère} & \textbf{Algo. glouton} & \textbf{Gale-Shapley} \\
\hline
Stabilité des affectations & \xmark & \cmark \\
\hline
Reproductibilité & \xmark & \cmark \\
\hline
Indépendance à l'ordre de traitement & \xmark & \cmark \\
\hline
Optimalité pour les étudiants & Résultat non garanti selon l'ordre de traitement & \cmark \\
\hline
Explicabilité des résultats & Difficile à justifier si l'ordre a influencé le résultat & \cmark \\
\hline
Simplicité d'implémentation & \cmark & Plus complexe que l'algo glouton, bien documenté \\
\hline
\end{tabular}
\end{table}

\subsection{Décision}

\textbf{Gale-Shapley} est retenu, dans sa variante \textbf{Student-Proposing} :
ce sont les étudiants qui font des propositions aux destinations, et non
l'inverse. Cette variante garantit que chaque étudiant obtient la meilleure
destination possible compte tenu de son classement et de ses vœux, ce qui
constitue la propriété dite d'\textit{optimalité côté proposant}.

Les contraintes de quota sont vérifiées à deux niveaux simultanément à
chaque proposition : un quota global par accord bilatéral (capacité de
l'université partenaire) et un quota local par département pédagogique
(places allouées au sein de l'ENSEEIHT). Une destination est disponible
seulement si les deux niveaux sont respectés ; sinon la proposition est
rejetée et l'étudiant passe à son vœu suivant.

L'algorithme glouton est écarté car il est sensible à l'ordre de traitement
et ne garantit pas la reproductibilité. Gale-Shapley assure un résultat
stable, équitable et reproductible dans ce cadre multi-contraint.


% ─────────────────────────────────────────────────────────────
\section{Choix de la méthode de répartition des quotas par département}
% ─────────────────────────────────────────────────────────────

Cette section présente les besoins liés à la répartition des quotas entre
départements, évalue les méthodes disponibles, puis justifie le choix retenu.

\subsection{Besoins spécifiques}

Chaque accord bilatéral fixe un nombre total de places disponibles. Ce nombre
doit ensuite être distribué entre les départements pédagogiques, ce qui soulève
une question qui n'a pas de réponse évidente : sur quelle base décider combien
de places revient à chaque département ?

Une répartition basée uniquement sur les demandes de l'année en cours serait
insuffisante : elle ne tient pas compte du fait qu'un département peut avoir été
systématiquement défavorisé les années précédentes, ni du fait que certains
accords voient leur nombre de places évoluer d'une campagne à l'autre. Sans
intégrer cet historique, la répartition risque d'être perçue comme arbitraire
et de creuser des déséquilibres entre départements au fil des années.

Le besoin est donc de disposer d'une méthode qui distribue les places de façon
proportionnelle à la demande réelle, tout en restant simple à comprendre,
transparente et vérifiable par le service RI. La méthode retenue doit également
garantir que la somme des places attribuées aux départements est exactement égale
au nombre de places de l'accord, sans qu'aucune place ne soit perdue lors des
arrondis.

\subsection{Alternatives évaluées}

Trois méthodes de répartition proportionnelle ont été comparées. Toutes
partent du même principe : chaque département reçoit une quote-part
proportionnelle à son nombre d'étudiants. La différence porte sur la
manière de distribuer les places entières lorsque la division produit
des résultats décimaux.

\subsubsection{Méthode Hamilton (plus grands restes)}

La méthode Hamilton attribue d'abord à chaque département la partie entière
de sa quote-part. Les places restantes sont ensuite données, une par une,
aux départements ayant les fractions les plus élevées, jusqu'à épuiser le
total disponible.

\textit{Avantages :}
\begin{itemize}
  \item Simple à comprendre et à expliquer à tous les acteurs
  \item Résultat facile à vérifier manuellement
  \item Respecte les proportions au plus près du total disponible
\end{itemize}

\textit{Limites :}
\begin{itemize}
  \item Un département dont la quote-part est inférieure à 1 ne reçoit
        aucune place garantie, même si sa fraction est élevée
\end{itemize}

\subsubsection{Méthode Jefferson (plus grandes moyennes)}

La méthode Jefferson procède de façon itérative : à chaque étape, la
prochaine place est attribuée au département dont le rapport entre sa
taille et le nombre de places déjà reçues est le plus élevé.

\textit{Avantages :}
\begin{itemize}
  \item Aucun arrondi explicite nécessaire
  \item Résultats stables et cohérents d'une campagne à l'autre
\end{itemize}

\textit{Limites :}
\begin{itemize}
  \item Avantage les grands départements et pénalise les petits
  \item Calcul itératif moins intuitif à vérifier manuellement
\end{itemize}

\subsubsection{Méthode Webster}

La méthode Webster est proche de Jefferson, mais elle utilise un arrondi
standard (seuil à 0,5) pour décider si un département reçoit une place
supplémentaire ou non.

\textit{Avantages :}
\begin{itemize}
  \item Plus équilibrée que Jefferson pour les petits départements
  \item Bien adaptée aux cas où peu de places sont à distribuer
\end{itemize}

\textit{Limites :}
\begin{itemize}
  \item Calcul moins intuitif qu'Hamilton, plus difficile à vérifier
        sans outil
  \item Résultats moins prévisibles dans les cas limites
\end{itemize}

\subsection{Synthèse comparative}

Le tableau~\ref{tab:comparaison_quotas} récapitule les critères évalués pour chacune des
trois méthodes, afin de faciliter la comparaison et de justifier le choix retenu.

\renewcommand{\arraystretch}{1.35}
\begin{table}[H]
\centering
\caption{Comparaison des méthodes de répartition des quotas}
\label{tab:comparaison_quotas}
\footnotesize
\begin{tabular}{|p{5cm}|>{\centering\arraybackslash}p{2.8cm}|>{\centering\arraybackslash}p{2.8cm}|>{\centering\arraybackslash}p{2.8cm}|}
\hline
\rowcolor{gray!20}
\textbf{Critère} & \textbf{Hamilton} & \textbf{Jefferson} & \textbf{Webster} \\
\hline
Simplicité de calcul & \cmark & \xmark & \xmark \\
\hline
Vérification manuelle possible & \cmark & Calcul itératif difficile à suivre & Moins prévisible dans les cas limites \\
\hline
Équité entre départements & \cmark & Favorise les grands départements & \cmark \\
\hline
Respect du total disponible & \cmark & \cmark & \cmark \\
\hline
Adapté à un petit nombre de places & \cmark & \xmark & \cmark \\
\hline
Transparence pour les acteurs & \cmark & \xmark & \xmark \\
\hline
\end{tabular}
\end{table}

\subsection{Décision}

\textbf{La méthode Hamilton} est retenue. Elle offre le meilleur équilibre
entre équité, simplicité et transparence dans ce contexte. Les résultats
peuvent être vérifiés à la main par le service RI, ce qui est important
pour un processus qui touche directement les chances de départ des étudiants.

La méthode Jefferson est écartée car elle avantage systématiquement les
grands départements, ce qui irait à l'encontre de l'équité visée. Webster
est plus difficile à vérifier sans outil de calcul. Hamilton est donc la
méthode la plus adaptée pour une répartition claire et justifiable.

% ─────────────────────────────────────────────────────────────
\section{Choix du moteur de recommandation}
% ─────────────────────────────────────────────────────────────

Cette section présente les besoins du module de recommandation,
évalue les approches disponibles, puis justifie la méthode retenue.

\subsection{Besoins spécifiques}

Le moteur de recommandation a pour objectif d'aider les étudiants à
choisir leurs destinations avant la saisie des vœux dans MoveON. Il
ne prend pas de décision, mais propose un classement des destinations
selon leur pertinence pour le profil de l'étudiant. Les recommandations
doivent s'appuyer sur le profil académique et sur les données historiques
disponibles, tout en restant interprétables pour l'étudiant.

\subsection{Alternatives évaluées}

Trois grandes familles d'approches ont été évaluées pour répondre à
ce besoin.

\subsubsection{Approche basée sur des règles (rule-based)}

Cette approche repose sur des règles métier simples définies manuellement,
telles que filtrer les destinations selon le département et classer par
taux d'acceptation historique.

\textit{Avantages :}
\begin{itemize}
  \item Simple à mettre en place, sans données d'apprentissage
  \item Facile à comprendre et à expliquer
\end{itemize}

\textit{Limites :}
\begin{itemize}
  \item Peu flexible, ne s'adapte pas aux profils individuels
  \item Résultats souvent trop génériques
\end{itemize}

\subsubsection{Approche par similarité (collaborative filtering)}

Cette approche propose des destinations en se basant sur les choix
d'étudiants ayant des profils similaires dans les promotions précédentes.

\textit{Avantages :}
\begin{itemize}
  \item Prend en compte les expériences passées des promotions
  \item Permet des recommandations plus personnalisées
\end{itemize}

\textit{Limites :}
\begin{itemize}
  \item Nécessite un volume important de données historiques
  \item Moins efficace pour les nouvelles destinations
  \item Mise en œuvre plus complexe
\end{itemize}

\subsubsection{Approche par apprentissage supervisé}

Cette approche consiste à entraîner un modèle pour prédire une probabilité
de compatibilité entre un étudiant et une destination. Deux modèles ont
été comparés au sein de cette famille.

\paragraph{Régression logistique}

La régression logistique est un modèle statistique qui estime la
probabilité d'un résultat à partir de variables d'entrée numériques.
Elle est particulièrement adaptée aux volumes de données modérés~\cite{scikit_learn}.

\textit{Avantages :}
\begin{itemize}
  \item Modèle simple et rapide à entraîner
  \item Résultats faciles à interpréter et à expliquer
  \item Adapté à des volumes de données modérés
\end{itemize}

\textit{Limites :}
\begin{itemize}
  \item Capture difficilement des relations complexes entre variables
\end{itemize}

\paragraph{XGBoost}

XGBoost est un algorithme de gradient boosting qui construit un ensemble
d'arbres de décision de façon séquentielle pour améliorer la précision
prédictive~\cite{scikit_learn}.

\textit{Avantages :}
\begin{itemize}
  \item Capable de modéliser des relations complexes entre variables
  \item Très bonnes performances prédictives
\end{itemize}

\textit{Limites :}
\begin{itemize}
  \item Plus complexe à entraîner et à maintenir
  \item Résultats moins interprétables
  \item Risque de surapprentissage avec un volume de données limité
\end{itemize}

\subsection{Synthèse comparative}

Le tableau~\ref{tab:comparaison_reco} récapitule les critères évalués pour chacune des
approches, afin de faciliter la comparaison et de justifier le choix retenu.

\renewcommand{\arraystretch}{1.35}
\begin{table}[H]
\centering
\caption{Comparaison des approches de recommandation}
\label{tab:comparaison_reco}
\footnotesize
\begin{tabular}{|p{3.8cm}|>{\centering\arraybackslash}p{2cm}|>{\centering\arraybackslash}p{2cm}|>{\centering\arraybackslash}p{2.5cm}|>{\centering\arraybackslash}p{2cm}|}
\hline
\rowcolor{gray!20}
\textbf{Critère} & \textbf{Règles} & \textbf{Collab.} & \textbf{Régression log.} & \textbf{XGBoost} \\
\hline
Personnalisation & \xmark & \cmark & \cmark & \cmark \\
\hline
Volume de données requis & \cmark & \xmark & Efficace à partir de 500 observations & Risque de surapprentissage avec peu de données \\
\hline
Interprétabilité & \cmark & Difficile sans connaître les profils similaires & \cmark & \xmark \\
\hline
Risque de surapprentissage & \cmark & Dépend de la qualité des données historiques & \cmark & \xmark \\
\hline
Simplicité de maintenance & \cmark & \xmark & \cmark & \xmark \\
\hline
\end{tabular}
\end{table}

\subsection{Décision}

\textbf{L'apprentissage supervisé par régression logistique} est retenu.
Le modèle estime, pour chaque couple étudiant–destination, une probabilité
de compatibilité à partir des variables suivantes :

\renewcommand{\arraystretch}{1.3}
\begin{table}[H]
\centering
\caption{Variables d'entrée du moteur de recommandation}
\label{tab:features_reco}
\footnotesize
\begin{tabular}{|p{3.5cm}|p{8cm}|}
\hline
\rowcolor{gray!20}
\textbf{Variable} & \textbf{Description} \\
\hline
GPA de l'étudiant & Note académique normalisée, principal signal de classement \\
\hline
Département pédagogique & Filtrage sur les accords ouverts au département de l'étudiant \\
\hline
Destination & Identifiant de l'université partenaire \\
\hline
Historique des affectations & Taux de sélection par destination sur les trois dernières années \\
\hline
\end{tabular}
\end{table}

Les destinations sont classées par ordre de probabilité décroissante
afin de proposer une liste de recommandations personnalisées.

La régression logistique est retenue car elle offre un bon compromis
entre simplicité, interprétabilité et facilité de maintenance. Le volume
de données, environ 900 observations sur trois ans, ne justifie pas
XGBoost, dont les performances supplémentaires s'accompagnent d'un risque
accru de surapprentissage.

\textbf{Risque de démarrage à froid.} Certaines destinations rares ou
récentes disposent de peu d'observations historiques. Pour ces cas, le
modèle se rabat sur le GPA et le département comme seuls signaux, ce qui
reste cohérent avec les règles de sélection. Le module sera activé à
partir de la deuxième année de fonctionnement, une fois qu'un historique
minimal d'une campagne complète aura été constitué.


% ─────────────────────────────────────────────────────────────
\section{Choix de l'orchestration RAG}
% ─────────────────────────────────────────────────────────────

Cette section présente les besoins du module RAG et le modèle de langage
imposé, évalue les frameworks d'orchestration disponibles, puis justifie
le choix retenu.

\subsection{Besoins spécifiques}

Le composant RAG du projet est un chatbot destiné aux étudiants, capable
de répondre à des questions sur les procédures de mobilité à partir de
la documentation institutionnelle de l'ENSEEIHT : guides, critères
d'éligibilité, délais et démarches administratives. Le cas d'usage est
relativement simple : une base documentaire stable, des questions en
français ou en anglais, et des réponses ancrées dans les documents officiels.

\subsection{Modèle de langage}

Le modèle de langage retenu est \textbf{Mistral}~\cite{mistral_dev}
(\texttt{mistral-small-latest}), conformément aux exigences de l'ENSEEIHT.
Il est accessible via l'\textbf{API Mistral ILAS}, une infrastructure
nationale dédiée aux établissements de l'enseignement supérieur français,
hébergée sur des serveurs localisés en France.

Ce mode d'accès institutionnel répond directement aux exigences RGPD :
les données traitées restent dans l'infrastructure nationale et ne sont
pas transmises à des tiers commerciaux. Il est néanmoins impératif que
les données personnelles des étudiants ne soient jamais incluses dans
les requêtes envoyées au modèle : les questions au chatbot sont
transmises telles que saisies, sans enrichissement avec des données
de dossier.


\subsection{Alternatives évaluées}

Quatre approches ont été évaluées pour orchestrer le pipeline RAG
autour du modèle Mistral.

\subsubsection{LangChain}

LangChain est un framework d'orchestration IA généraliste, compatible
avec un large éventail de modèles et de bases vectorielles~\cite{langchain_dev}.

\textit{Avantages :}
\begin{itemize}
  \item Grande communauté et nombreuses ressources disponibles
  \item Intégration simple avec Mistral et ChromaDB
  \item Adapté aux pipelines IA complexes et multi-étapes
\end{itemize}

\textit{Limites :}
\begin{itemize}
  \item API qui évolue fréquemment, induisant des risques de rupture
  \item Nombreuses dépendances, surcoût pour un chatbot documentaire simple
\end{itemize}

\subsubsection{LlamaIndex}

LlamaIndex est un framework spécialisé dans la recherche et la génération
à partir de documents structurés~\cite{LlamaIndex_dev}.

\textit{Avantages :}
\begin{itemize}
  \item Conçu spécifiquement pour les applications RAG documentaires
  \item Bonne intégration avec Mistral et les bases vectorielles
  \item API plus stable, prise en main plus directe
\end{itemize}

\textit{Limites :}
\begin{itemize}
  \item Moins adapté aux architectures IA multi-agents très complexes
\end{itemize}

\subsubsection{Haystack}

Haystack est un framework orienté production pour les systèmes de
recherche et de génération augmentée~\cite{haystack_dev}.

\textit{Avantages :}
\begin{itemize}
  \item Architecture modulaire adaptée aux déploiements en production
  \item Compatible avec Mistral
\end{itemize}

\textit{Limites :}
\begin{itemize}
  \item Documentation moins accessible pour des cas d'usage simples
  \item Intégration avec Mistral moins directe que LlamaIndex
\end{itemize}

\subsubsection{Approche directe (sans framework)}

L'approche directe consiste à implémenter manuellement le pipeline RAG
via des appels à l'API Mistral, sans couche d'abstraction intermédiaire.

\textit{Avantages :}
\begin{itemize}
  \item Contrôle total sur chaque étape du pipeline
  \item Peu de dépendances externes
\end{itemize}

\textit{Limites :}
\begin{itemize}
  \item Nécessite de développer manuellement toute la logique RAG
        (chunking, indexation, retrieval, construction du prompt)
  \item Maintenance plus complexe à long terme
\end{itemize}

\subsection{Synthèse comparative}

Le tableau~\ref{tab:comparaison_rag} récapitule les critères évalués pour chacune des
quatre approches, afin de faciliter la comparaison et de justifier
le choix retenu.

\renewcommand{\arraystretch}{1.35}
\begin{table}[H]
\centering
\caption{Comparaison des frameworks d'orchestration RAG}
\label{tab:comparaison_rag}
\footnotesize
\begin{tabular}{|p{4cm}|>{\centering\arraybackslash}p{2.2cm}|>{\centering\arraybackslash}p{2.2cm}|>{\centering\arraybackslash}p{2.2cm}|>{\centering\arraybackslash}p{2cm}|}
\hline
\rowcolor{gray!20}
\textbf{Critère} & \textbf{LangChain} & \textbf{LlamaIndex} & \textbf{Haystack} & \textbf{Direct} \\
\hline
Spécialisé RAG & Généraliste, non optimisé pour le RAG & \cmark & Orienté pipelines complexes & \xmark \\
\hline
Intégration Mistral & \cmark & \cmark & Possible mais moins directe & \cmark \\
\hline
Stabilité de l'API & \xmark & \cmark & \cmark & \cmark \\
\hline
Adapté à un cas simple & Surcoût de dépendances pour un cas simple & \cmark & Trop complexe pour un chatbot documentaire & Logique RAG à développer entièrement \\
\hline
Maintenance à long terme & API instable, risque de rupture à chaque mise à jour & \cmark & \cmark & \xmark \\
\hline
\end{tabular}
\end{table}

\subsection{Décision}

\textbf{LlamaIndex} est retenu comme framework d'orchestration RAG,
avec \textbf{ChromaDB} comme base vectorielle et
\textbf{Mistral} comme modèle de langage. Cette combinaison correspond
directement au besoin du projet : construire un chatbot documentaire
simple, basé sur une base de connaissances stable. LlamaIndex fournit
les fonctionnalités nécessaires sans surcoût de complexité, avec une
API stable et une intégration directe avec Mistral et ChromaDB.

% ─────────────────────────────────────────────────────────────
\section{Synthèse des choix algorithmiques et de performance}
% ─────────────────────────────────────────────────────────────

\renewcommand{\arraystretch}{1.4}

\begin{longtable}{|p{3cm}|p{3.5cm}|p{8.5cm}|}
\caption{Synthèse des choix algorithmiques et de performance} \\
\hline
\rowcolor{gray!20}
\textbf{Composant} & \textbf{Méthode retenue} & \textbf{Justification principale} \\
\hline
\endfirsthead

\hline
\rowcolor{gray!20}
\textbf{Composant} & \textbf{Méthode retenue} & \textbf{Justification principale} \\
\hline
\endhead

Stratégie de performance BDD & Indexation ciblée & Solution simple et suffisante pour le volume de données du projet \\
\hline

Algorithme d'affectation & Gale-Shapley & Stabilité, équité et reproductibilité des affectations \\
\hline

Répartition des quotas & Hamilton & Méthode proportionnelle simple, transparente et vérifiable manuellement \\
\hline

Moteur de recommandation & Régression logistique & Modèle simple, interprétable et adapté à un volume de données limité \\
\hline

Orchestration RAG & LlamaIndex + ChromaDB & Solution adaptée à un chatbot documentaire simple et stable \\
\hline

Modèle de langage & Mistral & Compatible avec le besoin multilingue du chatbot étudiant \\
\hline

\end{longtable}

\section*{Conclusion}
\justifying
\phantomsection
\addcontentsline{toc}{section}{Conclusion}

Ce chapitre a justifié l'ensemble des choix technologiques et algorithmiques qui structurent la plateforme, en privilégiant la cohérence de l'ensemble, l'équité et la simplicité d'exploitation. Ces choix forment le socle sur lequel repose la conception détaillée du système, présentée au chapitre suivant.
